\begin{frame}{자주 묻는 질문 --- PCA (2/2)}
  \small
  \begin{itemize}
    \item \kb{Q. 정보는 완전히 보존되지 않는 건가?}
      PCA는 \kb{손실 있는 압축}(lossy) --- 상위 $d$개만 남기면 나머지는
      \textbf{영구히} 사라짐. 그러나 ``분산 기준''으로 중요 부분부터
      보존하므로, 잃는 양은 \textbf{버린 고유값의 합}으로 정확히 계산
      가능 --- 무작위로 잘라내는 것과 달리 ``얼마나 잃는지''가 정량적.
      \par \textbf{모든} 정보를 보존하려면(lossless) 전체 주성분을
      모두 남겨야 한다.
    \item \kb{Q. k-means 전에 항상 필요한가?}
      아니다. 특히 잘 맞는 경우: (a) 특징끼리 \textbf{상관}되어 중복이
      있을 때, (b) 차원이 높아 \textbf{차원의 저주}가 문제일 때,
      (c) \textbf{노이즈}(분산이 작은 성분) 제거가 필요할 때.
      \par 특징이 서로 무관하고 차원이 낮다면 불필요한 단계(오히려
      해석성을 잃을 수 있음).
  \end{itemize}
\end{frame}
